
The most general approach to modeling warping is to take the \(n_w\) warping amplitudes \(\bm{\alpha}\) as independent configuration variables. 
This generalizes the approach of \cite{reissner1952nonuniform} for warping torsion.

\paragraph{Equilibrium}
Regardless of any kinematic assumptions on the primary kinematic variables \(\bm{x}\) and \(\FrameRotation\), the warping resultants \(\bm{w}\) and \(\bm{v}\) are governed by the same equilibrium equation:
\begin{equation}
\bm{w}' - \bm{v} + \bm{f}_{\mathrm{w}} = \mathbf{0}
\label{eq:bi-equilibrium}
\end{equation}
where $\bm{f}_{\mathrm{w}}$ is a generalized distributed load in the warping variables (typically taken to be zero). 
%
The complete system of equilibrium equations is
\begin{equation} 
\left\{\begin{array}{r}
(\FrameRotation \bm{n})^{\prime}+{\bm{f}}_n=\mathbf{0} \\
(\FrameRotation \bm{m})^{\prime}+\bm{x}^{\prime} \times (\FrameRotation \bm{n}) +{\bm{f}}_m=\mathbf{0} \\
\bm{w}^{\prime}-\bm{v} + \bm{f}_{\rm w} = \mathbf{0}
\end{array}\right. 
\label{eq:finite-equilibrium}
\end{equation}
where \(\bm{f}_n\) and \(\bm{f}_m\) are given body force and moment loadings, respectively. 
\citep{antman1974kirchhoffs,simo1991geometricallyexact,klinkel2003anisotropic}. 

% This is separated out of the equilibrium equation \eqref{eq:finite-equilibrium}, so that 
% all assumptions on and simplifications to the primary kinematic variables \(\bm{x}\) and \(\FrameRotation\) are contained within a \emph{primary} equilibrium operator \(\mathbb{B}_{\mathrm{p}}\). 


\vspace{1ex}
\begin{ambox}[Type 1: Independent Warping]{box:type1-warping}
A Type 1 constitutive response maps \(6+2n_w\) strains \((\bm{\gamma},\, \bm{\kappa},\, \bm{\alpha}',\, \bm{\alpha})\) to \(6+ 2n_w\) generalized stresses \((\bm{n},\, \bm{m},\,\bm{w},\bm{v})\)
\[
[\bm{n},\bm{m},\bm{w},\bm{v}] = \mathcal{S}(\bm{\gamma},\, \bm{\kappa},\, \bm{\alpha}',\, \bm{\alpha})
\]

\begin{equation}
\begin{aligned}
    \mathbb{B}_{\mathrm{p}} \bm{A}_*\begin{pmatrix} \bm{n} \\ \bm{m} \end{pmatrix} + \bm{f}_{\mathrm{p}} &= \dot{\bm{\mu}} \\ 
    \bm{w}' - \bm{v} + \bm{f}_{\mathrm{w}} &= \mathbf{0}
\end{aligned}
\end{equation}
\end{ambox}


\subsection{Type 1: Independent Warping}

The most general approach to model warping is to take the \(\nwarp\) warping amplitudes \(\bm{\alpha}\) as independent configuration variables. 
This follows the approach of \citep{reissner1952nonuniform} for warping torsion. 
Regardless of any kinematic assumptions on the primary kinematic variables \(\TraceLocation\) and \(\FrameRotation\), the warping resultants \(\bimoment\) and \(\bishear\) are governed by the same equilibrium equation:
\begin{equation*}
    \bimoment' - \bishear + \bm{f}_{\mathrm{w}} = \mathbf{0}
\end{equation*}
This is separated out of the equilibrium equation \eqref{eq:finite-equilibrium}, so that all assumptions on and simplifications to the primary kinematic variables \(\TraceLocation\) and \(\FrameRotation\) are contained within a \emph{primary} equilibrium operator \(\EquilibriumDerivative\). 

A \textbf{Type 1} constitutive formulation is completely general in the sense that it can readily return a response for all other formulation types.

\vspace{1ex}
\begin{ambox}[Type 1: Independent Warping]{box:type1-warping}
A Type 1 constitutive response maps \(6+2\, \nwarp\) strains \((\ShearStrain,\, \CurveStrain,\, \bm{\alpha}',\, \bm{\alpha})\) to \(6+ 2\, \nwarp\) generalized stresses \((\ShearResult,\, \CurveResult,\,\bm{w},\bm{v})\)
\[
[\bm{s}_{\mathrm{p}},\bm{w},\bm{v}] = \SectionEvaluation{\bm{e}_{\mathrm{p}},\, \bm{\alpha}',\, \bm{\alpha}}
\]
Equilibrium takes the form 
\begin{equation}
\begin{aligned}
    \mathbb{B}_{\rm p} \bm{\sigma}_{\mathrm{p}} + \bm{f}_{\mathrm{p}} &= \dot{\bm{\mu}} \\ 
    \bimoment' - \bishear + \bm{f}_{\mathrm{w}} &= \mathbf{0}
\end{aligned}
\end{equation}
\end{ambox}

